\documentclass[10pt,a4paper]{amsart}
\usepackage[margin=19mm]{geometry}
\usepackage{amsmath,amssymb,mathtools}
\usepackage[scr=rsfs,cal=euler]{mathalfa}
\usepackage{xfrac}
\usepackage[unicode,hidelinks,bookmarks=false]{hyperref}
\newcommand{\ie}{i.e.\ }
\newcommand{\cf}{cf.\ }
\newcommand{\resp}{respectively\ }
\newcommand{\Subsection}{\subsection}
\providecommand{\nobreakdash}{\nobreak\hbox{-}}
\setlength{\emergencystretch}{2em}
\multlinegap0pt
\usepackage{bm}
\usepackage{bbm}
\usepackage{dsfont}
\usepackage{xspace}
\usepackage{cancel}
\usepackage{mathtools}
\usepackage{amsthm}
\makeatletter
\@ifundefined{theorem}{\newtheorem{theorem}{Theorem}}{}
\@ifundefined{lemma}{\newtheorem{lemma}[theorem]{Lemma}}{}
\makeatother
\DeclareMathOperator{\Tr}{Tr}
\title[Geodesics in the Deep Linear Network]{Geodesics in the Deep Linear Network}
\author{Alan Chen}
\date{Source: arXiv:2510.07324v2 (CC BY 4.0)}
\begin{document}
\maketitle

\noindent\emph{Source and licence.} Geodesics in the Deep Linear Network. Version arXiv:2510.07324v2 (CC BY 4.0), distributed under the Creative Commons Attribution 4.0 International licence, as recorded in the arXiv licence metadata for this exact version and shown on the abstract page https://arxiv.org/abs/2510.07324v2. Full rights notice: licenses/arxiv-2510.07324-CC-BY-4.0.txt.\par\smallskip

\noindent\textit{Editorial scope.} Counted unit: the Theorem whose source label is \texttt{thm:bw\_general\_case} in the article (source0.tex lines 279--293, source label \texttt{thm:bw\_general\_case}), delivered together with the article's own complete original proof of it (source0.tex lines 295--348, one \texttt{proof} environment of 54 lines). No mathematics is written, repaired or re-derived: statement and proof are the article's own text line for line. Required context retained uncounted: lem:geometric\_mean\_riccati (lines 269--277). Every definition, display and label of those lines is retained unchanged. Omitted from this package and not redistributed: 1--268, 278--278, 294--294, 349--646. Nothing inside a retained excerpt is omitted or altered: every retained body is delivered exactly as the article's lines are written, so no omission receipt of any kind accompanies this package. Citations: the retained excerpts carry no citation command at all, so no bibliography entry of the article is transcribed and no reference is redistributed. Reference adaptation: none was needed and none was made; every reference inside the delivered unit resolves inside it, so no unresolved reference or citation is produced. Printed slips kept literal: no printed word, symbol, hypothesis or number of the counted statement or of its proof was altered. Layout and class: the article's own class is \texttt{amsart} with packages including xcolor; the wrapper is the lane's pinned \texttt{amsart} editorial wrapper at 10pt on A4, which restates the theorem-like environments the retained text needs and adds the article's own macro definitions that the retained text needs (\texttt{\textbackslash Tr}), with extra wrapper packages (bm, bbm, dsfont, xspace, cancel, mathtools). The delivered numbering is the unit's own. Assets: no retained span contains a figure environment, a graphics-inclusion command, a PDF-page inclusion, a table or a TikZ picture; no image file is copied into this package, the package declares no asset dependency and no image file is redistributed with it. Licence: version arXiv:2510.07324v2 of this article is distributed under the Creative Commons Attribution 4.0 International licence, as recorded in the arXiv licence metadata for this exact version and shown on the abstract page https://arxiv.org/abs/2510.07324v2. A full rights notice is delivered at licenses/arxiv-2510.07324-CC-BY-4.0.txt. Characters: every retained excerpt is pure ASCII, so the delivered engine, XeLaTeX with its default Latin Modern text font, reports no missing character.
\begin{lemma}\label{lem:geometric_mean_riccati}
\begin{enumerate}
    \item $A\#B$ is the unique positive definite solution $M$ to the Riccati Equation
    \begin{equation}
        B = MA^{-1}M.
    \end{equation}
    \item $A(A^{-1}\#B) = (AB)^{1/2}$ and $(A^{-1}\#B)A = (BA)^{1/2}$.
\end{enumerate}
\end{lemma}

\begin{theorem}[Bures-Wasserstein Geodesics]\label{thm:bw_general_case}
    The geodesic path $X: [0,1] \to \mathbb{P}_d$ between $A, B \in \mathbb{P}_d$ in $(\mathbb{P}_d, g^\text{BW})$ satisfies the boundary value problem
    \begin{equation}\label{eq:bures_wasserstein_equations}
        \begin{cases}
            \dot{X} = \mathcal{L}_X(P), \\
            \dot{P} = -P^2,\\
            X(0) = A,\\
            X(1) = B,
        \end{cases}
    \end{equation}
    and has explicit form given by
\begin{equation}
    X(t) = (1-t)^2 A + t^2 B + t(1-t)\left((AB)^{1/2} + (BA)^{1/2}\right).
\end{equation}
\end{theorem}

\begin{proof}[Proof of Theorem \ref{thm:bw_general_case}]
To begin, recall that $T_X\mathbb{P}_d \cong \text{Symm}_d$ and denote $Y = \dot{X}$. We consider the Hamiltonian system, which from the inverse metric tensor $\mathcal{L}$ is
\begin{align}
    H(X, P) &= \frac{1}{2}\Tr(P\mathcal{L}_X(P)) = \Tr(P^2X).
\end{align}
The system of differential equations is now clear from the form of $H(X,P)$. We can solve them to find exact formulas for $X$ and $P$. In particular, $\dot{P} = -P^2$ implies
\begin{align}
    -P^{-1} \dot{P} P^{-1} &= I \nonumber \\
    \frac{d}{dt}(P^{-1}) &= I \nonumber \\
    P^{-1}(t) &= tI + P^{-1}_0.
\end{align}
Thus, $P(t) = (tI + P^{-1}_0)^{-1}$. Factoring out $P^{-1}_0$ from the interior term reveals that $P(t)$ is a time derivative:
\begin{align}
    P(t) &= (P^{-1}_0(tP_0 + I))^{-1}  \nonumber \\
    &= (tP_0 + I)^{-1}P_0 \\
    &= \frac{d}{dt} \log (tP_0 + I).
\end{align}
The first differential equation can be expanded as
\begin{align}
    \dot{X} = X(t)P(t) + P(t)X(t).
\end{align}
Motivated by the similarity of this equation to the Lax equation, we make the ansatz for the solution form of $X$ as
\begin{equation}
    X(t) = e^{C(t)} X_0 e^{D(t)},
\end{equation}
which after plugging in and matching coefficients gives that $\dot{C}(t) = \dot{D}(t) = P(t)$ and implying that 
\begin{align}
    C(t) = D(t) &= \int_0^t P(s) ds\\
    &= \int_{0}^t \frac{d}{ds}\log(sP_0 + I) ds \\
    &= \log (tP_0+I).
\end{align}
Substituting $C$ and $D$ back into the ansatz gives that 
\begin{align}
    X(t) &= \exp\{\log (tP_0+I)\} X_0 \exp\{\log (tP_0+I)\}\nonumber \\
    &= (tP_0 + I) X_0 (tP_0 + I) \nonumber \\ 
    &= (tP_0 + I) A (tP_0 + I).
\end{align}
Applying the final unused condition that $X_1 = B$ helps find $P_0$:
\begin{equation}
    B = (P_0 + I) A (P_0 + I),
\end{equation}
which from Lemma \ref{lem:geometric_mean_riccati} implies that
\begin{equation}
    P_0 = (A^{-1} \# B) - I.
\end{equation}
Finally, substituting $P_0$ into our solution formula for $X(t)$, we conclude that
\begin{align}
    X(t) &= (t((A^{-1} \# B) - I) + I) A (t((A^{-1} \# B) - I) + I) \\
    &= (t(A^{-1} \# B) + (1-t)I) A (t(A^{-1} \# B) + (1-t)I)\\
    &= (1-t)^2 A + t^2 B + t(1-t)((A^{-1}\#B)A + A(A^{-1}\#B))\\
    &= (1-t)^2 A + t^2 B + t(1-t)\left((AB)^{1/2} + (BA)^{1/2}\right).
\end{align}
The uniqueness of this solution also follows from Picard's theorem, as the right hand sides of Equation \ref{eq:bures_wasserstein_equations} are smooth.
\end{proof}

\end{document}
